% explainer-source-hash: 30eba2c07fd93bd3e818b2a9942e6f70ace175b93b997cfd1d710afb07bbfc59
% explainer-source-commit: 6a20de2f6ebc03e4d85c9454ba07bb53ed18f9da
% explainer-generated: 2026-10-15
% Regenerate with /explain-all (see WIRING.md). Do not edit this stamp by hand:
% it is how the toolchain knows whether this document still matches the code.
\documentclass[11pt,a4paper]{article}
% preamble.tex - shared preamble for this project's explainer documents.
%
% This file is generated: it is copied in beside the documents on every build, so a
% clone of this repo can rebuild the PDFs, but local edits to it are overwritten.
%
% Usage from a document:
%   \documentclass[11pt,a4paper]{article}
%   % preamble.tex - shared preamble for this project's explainer documents.
%
% This file is generated: it is copied in beside the documents on every build, so a
% clone of this repo can rebuild the PDFs, but local edits to it are overwritten.
%
% Usage from a document:
%   \documentclass[11pt,a4paper]{article}
%   % preamble.tex - shared preamble for this project's explainer documents.
%
% This file is generated: it is copied in beside the documents on every build, so a
% clone of this repo can rebuild the PDFs, but local edits to it are overwritten.
%
% Usage from a document:
%   \documentclass[11pt,a4paper]{article}
%   \input{preamble}
%   \begin{document} ... \end{document}
%
% Build:  pdflatex -interaction=nonstopmode deep-dive.tex   (run it three times, so the
%         table of contents and the TikZ positioning settle)
%
% PACKAGE NOTES. Every package used here was checked present on the machine these
% documents were built on (Debian TeX Live 2023). If you rebuild elsewhere and a build
% fails, these are the likely gaps. Do not add a package without checking first:
%   kpsewhich <pkg>.sty
%
% Deliberately NOT used, because they were absent and the build dies without them:
%   algorithm2e, algorithm, algorithmic, algpseudocode  (texlive-science not installed)
%     -> the \begin{pseudocode} environment defined below stands in for them.
%   lmodern, charter, mathptmx, times, mathpazo, courier, newtx*, tgtermes
%     -> their .sty files RESOLVE but the font files themselves are absent, so the
%        build dies with "I can't find file `bchr8t'" or similar. Verified by
%        test-compiling each one. Only Computer Modern (the default) and helvet
%        actually render. kpsewhich reporting a font it cannot render is the trap:
%        do not "upgrade" the font stack without test-compiling it first.
%   amsmath, amssymb -> not loaded, so \text{} and \blacksquare are unavailable.
%
% A fuller font stack and algorithm2e would come from:
%   sudo apt install texlive-fonts-recommended texlive-science cm-super

\usepackage[T1]{fontenc}
\usepackage[utf8]{inputenc}
% Font: Computer Modern (default). See the note above before changing this.
\usepackage[a4paper,margin=2.4cm,headheight=14pt]{geometry}
% microtype WITHOUT font expansion. This is load-bearing, do not "restore" it:
% cm-super is not installed, so T1 Computer Modern resolves to BITMAP fonts, and
% pdfTeX's auto-expansion needs scalable ones. With expansion on, any document
% long enough to reach a second page dies with
%   "pdfTeX error (font expansion): auto expansion is only possible with scalable fonts"
% Protrusion (the bigger visual win) still works and stays on.
% `sudo apt install cm-super` would lift this; until then, leave expansion off.
\usepackage[protrusion=true,expansion=false]{microtype}
% Kill hyphen ligatures in TYPEWRITER text. This is load-bearing, not cosmetic. Under T1,
% \texttt{--verbose} silently renders as "-verbose", and \texttt{---} as "--", so a
% document can print a command-line flag as the WRONG flag and look perfectly fine doing
% it. Verified by rendering the page and reading the glyphs. lstlisting and \verb were
% always safe; this makes \code, \file and \texttt safe too, which is what prose uses.
% Ordinary prose is unaffected and still ligates, which is what you want there.
\DisableLigatures[-]{family=tt*}
\usepackage{graphicx}
\usepackage{xcolor}
\usepackage{booktabs}
\usepackage{enumitem}
\usepackage{fancyhdr}
\usepackage{titlesec}
\usepackage{tcolorbox}
\usepackage{listings}
\usepackage{float}
\usepackage{caption}
\usepackage{tikz}
\usepackage{forest}
\usepackage{pgfplots}
\pgfplotsset{compat=1.18}
\usetikzlibrary{arrows.meta,positioning,shapes.geometric,fit,backgrounds,calc,decorations.pathreplacing}
\tcbuselibrary{skins,breakable}
\usepackage[hidelinks,bookmarks=true,bookmarksnumbered=true]{hyperref}

% ===== palette ===========================================================
\definecolor{inkblue}{HTML}{1F3A5F}
\definecolor{accent}{HTML}{2E7DAF}
\definecolor{softgrey}{HTML}{F4F6F8}
\definecolor{linegrey}{HTML}{C8D0D8}
\definecolor{codebg}{HTML}{F7F7F4}
\definecolor{warnred}{HTML}{A33B3B}
\definecolor{okgreen}{HTML}{2E7D53}

% ===== headings ==========================================================
\titleformat{\section}{\Large\bfseries\color{inkblue}}{\thesection}{0.7em}{}
\titleformat{\subsection}{\large\bfseries\color{inkblue}}{\thesubsection}{0.6em}{}
\titleformat{\subsubsection}{\normalsize\bfseries\color{accent}}{\thesubsubsection}{0.5em}{}

\pagestyle{fancy}
\fancyhf{}
\fancyhead[L]{\small\color{inkblue}\nouppercase{\leftmark}}
\fancyhead[R]{\small\color{inkblue}\thepage}
\renewcommand{\headrulewidth}{0.4pt}
\renewcommand{\headrule}{\hbox to\headwidth{\color{linegrey}\leaders\hrule height \headrulewidth\hfill}}

% ===== code ==============================================================
\lstdefinestyle{house}{
  backgroundcolor=\color{codebg},
  basicstyle=\ttfamily\footnotesize,
  keywordstyle=\color{inkblue}\bfseries,
  commentstyle=\color{gray}\itshape,
  stringstyle=\color{okgreen},
  numbers=left, numberstyle=\tiny\color{gray}, numbersep=8pt,
  frame=single, rulecolor=\color{linegrey},
  breaklines=true, breakatwhitespace=false,
  showstringspaces=false, tabsize=2,
  captionpos=b, xleftmargin=14pt, framexleftmargin=14pt,
  columns=fullflexible, keepspaces=true,
  upquote=true,  % render ' and " as straight quotes, not curly
  % NON-ASCII INSIDE A LISTING IS A HARD BUILD FAILURE, not a cosmetic issue: listings
  % reads raw bytes, so inputenc dies on the first multi-byte character ("Invalid UTF-8
  % byte sequence") and the whole document fails to build. These mappings cover what
  % realistically turns up in source code: Latin-1 accents, Spanish punctuation, smart
  % quotes picked up from a word processor, the degree and section signs.
  % NOT COVERED, and it cannot be here: Urdu, Arabic and Balochi script. pdflatex has no
  % font for them. Describe or transliterate such text; never paste it into a listing.
  literate=%
    {á}{{\'a}}1 {é}{{\'e}}1 {í}{{\'i}}1 {ó}{{\'o}}1 {ú}{{\'u}}1
    {Á}{{\'A}}1 {É}{{\'E}}1 {Í}{{\'I}}1 {Ó}{{\'O}}1 {Ú}{{\'U}}1
    {à}{{\`a}}1 {è}{{\`e}}1 {ì}{{\`i}}1 {ò}{{\`o}}1 {ù}{{\`u}}1
    {ä}{{\"a}}1 {ë}{{\"e}}1 {ï}{{\"i}}1 {ö}{{\"o}}1 {ü}{{\"u}}1
    {Ä}{{\"A}}1 {Ö}{{\"O}}1 {Ü}{{\"U}}1 {ß}{{\ss}}1
    {ñ}{{\~n}}1 {Ñ}{{\~N}}1 {ç}{{\c c}}1 {Ç}{{\c C}}1
    {¿}{{\textquestiondown}}1 {¡}{{\textexclamdown}}1
    {°}{{\textdegree}}1 {§}{{\S}}1 {£}{{\pounds}}1 {€}{{EUR}}1
    {“}{{\textquotedblleft}}1 {”}{{\textquotedblright}}1
    {‘}{{\textquoteleft}}1 {’}{{\textquoteright}}1
    {–}{{-}}1 {—}{{-}}1 {…}{{...}}1 {→}{{$\rightarrow$}}1 {✓}{{x}}1
}
\lstset{style=house}

% ===== callout boxes =====================================================
% \begin{keyidea}{Title} ... \end{keyidea}   - a concept a beginner must grasp
\newtcolorbox{keyidea}[1]{
  enhanced, breakable, colback=softgrey, colframe=accent,
  boxrule=0.4pt, arc=2pt, left=8pt, right=8pt, top=6pt, bottom=6pt,
  title={\bfseries #1}, coltitle=white, colbacktitle=accent, fonttitle=\small
}

% \begin{jargon}{Term} definition \end{jargon}  - define a term at first use
\newtcolorbox{jargon}[1]{
  enhanced, breakable, colback=white, colframe=linegrey,
  boxrule=0.4pt, arc=2pt, left=8pt, right=8pt, top=6pt, bottom=6pt,
  title={\bfseries Jargon: #1}, coltitle=inkblue, colbacktitle=softgrey, fonttitle=\small
}

% \begin{honest}{Title} ... \end{honest}  - brutally honest finding (rule 18)
\newtcolorbox{honest}[1]{
  enhanced, breakable, colback=white, colframe=warnred,
  boxrule=0.6pt, arc=2pt, left=8pt, right=8pt, top=6pt, bottom=6pt,
  title={\bfseries #1}, coltitle=white, colbacktitle=warnred, fonttitle=\small
}

% \begin{pseudocode}{Caption} ... \end{pseudocode}
% Stands in for algorithm2e, which is NOT installed. Use \Step / \Indent inside.
\newtcolorbox{pseudocode}[1]{
  enhanced, breakable, colback=codebg, colframe=inkblue,
  boxrule=0.4pt, arc=2pt, left=10pt, right=8pt, top=6pt, bottom=6pt,
  fontupper=\ttfamily\small,
  title={\bfseries\rmfamily #1}, coltitle=white, colbacktitle=inkblue, fonttitle=\small
}
\newcommand{\Step}[1]{\par\noindent #1}
\newcommand{\Indent}[1]{\par\noindent\hspace*{2em}#1}

% ===== tikz house styles =================================================
\tikzset{
  box/.style   = {rectangle, rounded corners=2pt, draw=inkblue, fill=softgrey,
                  text=inkblue, align=center, inner sep=6pt, minimum height=9mm,
                  font=\small},
  extbox/.style= {rectangle, rounded corners=2pt, draw=linegrey, fill=white,
                  text=black!70, align=center, inner sep=6pt, minimum height=9mm,
                  font=\small, dashed},
  store/.style = {cylinder, shape border rotate=90, aspect=0.25, draw=accent,
                  fill=white, text=inkblue, align=center, inner sep=5pt,
                  font=\small},
  dec/.style   = {diamond, aspect=1.6, draw=accent, fill=white, text=inkblue,
                  align=center, inner sep=2pt, font=\scriptsize},
  flow/.style  = {-{Stealth[length=2.4mm]}, draw=inkblue, thick},
  dflow/.style = {-{Stealth[length=2.4mm]}, draw=accent, thick, dashed},
  % align=center is load-bearing: without it, a \\ inside an edge label throws
  % the misleading "perhaps a missing \item" error.
  lbl/.style   = {font=\scriptsize\itshape, text=black!65, fill=white,
                  inner sep=2pt, align=center}
}

% forest default for directory trees
\forestset{
  dirtree/.style={
    for tree={
      font=\ttfamily\small, grow'=0, child anchor=west, parent anchor=south,
      anchor=west, calign=first, inner sep=2pt,
      edge path={\noexpand\path [draw, \forestoption{edge}]
        (!u.south west) +(9pt,0) |- (.child anchor) \forestoption{edge label};},
      before typesetting nodes={if n=1{insert before={[,phantom]}}{}},
      fit=band, before computing xy={l=15pt}
    }
  }
}

% ===== misc ==============================================================
\captionsetup{font=small, labelfont={bf,color=inkblue}}
\setlist[itemize]{leftmargin=1.2em, itemsep=2pt, topsep=3pt}
\setlist[enumerate]{leftmargin=1.4em, itemsep=2pt, topsep=3pt}
\newcommand{\file}[1]{\texttt{\small #1}}
\newcommand{\code}[1]{\texttt{\small #1}}
\setcounter{tocdepth}{2}
\setcounter{secnumdepth}{3}

%   \begin{document} ... \end{document}
%
% Build:  pdflatex -interaction=nonstopmode deep-dive.tex   (run it three times, so the
%         table of contents and the TikZ positioning settle)
%
% PACKAGE NOTES. Every package used here was checked present on the machine these
% documents were built on (Debian TeX Live 2023). If you rebuild elsewhere and a build
% fails, these are the likely gaps. Do not add a package without checking first:
%   kpsewhich <pkg>.sty
%
% Deliberately NOT used, because they were absent and the build dies without them:
%   algorithm2e, algorithm, algorithmic, algpseudocode  (texlive-science not installed)
%     -> the \begin{pseudocode} environment defined below stands in for them.
%   lmodern, charter, mathptmx, times, mathpazo, courier, newtx*, tgtermes
%     -> their .sty files RESOLVE but the font files themselves are absent, so the
%        build dies with "I can't find file `bchr8t'" or similar. Verified by
%        test-compiling each one. Only Computer Modern (the default) and helvet
%        actually render. kpsewhich reporting a font it cannot render is the trap:
%        do not "upgrade" the font stack without test-compiling it first.
%   amsmath, amssymb -> not loaded, so \text{} and \blacksquare are unavailable.
%
% A fuller font stack and algorithm2e would come from:
%   sudo apt install texlive-fonts-recommended texlive-science cm-super

\usepackage[T1]{fontenc}
\usepackage[utf8]{inputenc}
% Font: Computer Modern (default). See the note above before changing this.
\usepackage[a4paper,margin=2.4cm,headheight=14pt]{geometry}
% microtype WITHOUT font expansion. This is load-bearing, do not "restore" it:
% cm-super is not installed, so T1 Computer Modern resolves to BITMAP fonts, and
% pdfTeX's auto-expansion needs scalable ones. With expansion on, any document
% long enough to reach a second page dies with
%   "pdfTeX error (font expansion): auto expansion is only possible with scalable fonts"
% Protrusion (the bigger visual win) still works and stays on.
% `sudo apt install cm-super` would lift this; until then, leave expansion off.
\usepackage[protrusion=true,expansion=false]{microtype}
% Kill hyphen ligatures in TYPEWRITER text. This is load-bearing, not cosmetic. Under T1,
% \texttt{--verbose} silently renders as "-verbose", and \texttt{---} as "--", so a
% document can print a command-line flag as the WRONG flag and look perfectly fine doing
% it. Verified by rendering the page and reading the glyphs. lstlisting and \verb were
% always safe; this makes \code, \file and \texttt safe too, which is what prose uses.
% Ordinary prose is unaffected and still ligates, which is what you want there.
\DisableLigatures[-]{family=tt*}
\usepackage{graphicx}
\usepackage{xcolor}
\usepackage{booktabs}
\usepackage{enumitem}
\usepackage{fancyhdr}
\usepackage{titlesec}
\usepackage{tcolorbox}
\usepackage{listings}
\usepackage{float}
\usepackage{caption}
\usepackage{tikz}
\usepackage{forest}
\usepackage{pgfplots}
\pgfplotsset{compat=1.18}
\usetikzlibrary{arrows.meta,positioning,shapes.geometric,fit,backgrounds,calc,decorations.pathreplacing}
\tcbuselibrary{skins,breakable}
\usepackage[hidelinks,bookmarks=true,bookmarksnumbered=true]{hyperref}

% ===== palette ===========================================================
\definecolor{inkblue}{HTML}{1F3A5F}
\definecolor{accent}{HTML}{2E7DAF}
\definecolor{softgrey}{HTML}{F4F6F8}
\definecolor{linegrey}{HTML}{C8D0D8}
\definecolor{codebg}{HTML}{F7F7F4}
\definecolor{warnred}{HTML}{A33B3B}
\definecolor{okgreen}{HTML}{2E7D53}

% ===== headings ==========================================================
\titleformat{\section}{\Large\bfseries\color{inkblue}}{\thesection}{0.7em}{}
\titleformat{\subsection}{\large\bfseries\color{inkblue}}{\thesubsection}{0.6em}{}
\titleformat{\subsubsection}{\normalsize\bfseries\color{accent}}{\thesubsubsection}{0.5em}{}

\pagestyle{fancy}
\fancyhf{}
\fancyhead[L]{\small\color{inkblue}\nouppercase{\leftmark}}
\fancyhead[R]{\small\color{inkblue}\thepage}
\renewcommand{\headrulewidth}{0.4pt}
\renewcommand{\headrule}{\hbox to\headwidth{\color{linegrey}\leaders\hrule height \headrulewidth\hfill}}

% ===== code ==============================================================
\lstdefinestyle{house}{
  backgroundcolor=\color{codebg},
  basicstyle=\ttfamily\footnotesize,
  keywordstyle=\color{inkblue}\bfseries,
  commentstyle=\color{gray}\itshape,
  stringstyle=\color{okgreen},
  numbers=left, numberstyle=\tiny\color{gray}, numbersep=8pt,
  frame=single, rulecolor=\color{linegrey},
  breaklines=true, breakatwhitespace=false,
  showstringspaces=false, tabsize=2,
  captionpos=b, xleftmargin=14pt, framexleftmargin=14pt,
  columns=fullflexible, keepspaces=true,
  upquote=true,  % render ' and " as straight quotes, not curly
  % NON-ASCII INSIDE A LISTING IS A HARD BUILD FAILURE, not a cosmetic issue: listings
  % reads raw bytes, so inputenc dies on the first multi-byte character ("Invalid UTF-8
  % byte sequence") and the whole document fails to build. These mappings cover what
  % realistically turns up in source code: Latin-1 accents, Spanish punctuation, smart
  % quotes picked up from a word processor, the degree and section signs.
  % NOT COVERED, and it cannot be here: Urdu, Arabic and Balochi script. pdflatex has no
  % font for them. Describe or transliterate such text; never paste it into a listing.
  literate=%
    {á}{{\'a}}1 {é}{{\'e}}1 {í}{{\'i}}1 {ó}{{\'o}}1 {ú}{{\'u}}1
    {Á}{{\'A}}1 {É}{{\'E}}1 {Í}{{\'I}}1 {Ó}{{\'O}}1 {Ú}{{\'U}}1
    {à}{{\`a}}1 {è}{{\`e}}1 {ì}{{\`i}}1 {ò}{{\`o}}1 {ù}{{\`u}}1
    {ä}{{\"a}}1 {ë}{{\"e}}1 {ï}{{\"i}}1 {ö}{{\"o}}1 {ü}{{\"u}}1
    {Ä}{{\"A}}1 {Ö}{{\"O}}1 {Ü}{{\"U}}1 {ß}{{\ss}}1
    {ñ}{{\~n}}1 {Ñ}{{\~N}}1 {ç}{{\c c}}1 {Ç}{{\c C}}1
    {¿}{{\textquestiondown}}1 {¡}{{\textexclamdown}}1
    {°}{{\textdegree}}1 {§}{{\S}}1 {£}{{\pounds}}1 {€}{{EUR}}1
    {“}{{\textquotedblleft}}1 {”}{{\textquotedblright}}1
    {‘}{{\textquoteleft}}1 {’}{{\textquoteright}}1
    {–}{{-}}1 {—}{{-}}1 {…}{{...}}1 {→}{{$\rightarrow$}}1 {✓}{{x}}1
}
\lstset{style=house}

% ===== callout boxes =====================================================
% \begin{keyidea}{Title} ... \end{keyidea}   - a concept a beginner must grasp
\newtcolorbox{keyidea}[1]{
  enhanced, breakable, colback=softgrey, colframe=accent,
  boxrule=0.4pt, arc=2pt, left=8pt, right=8pt, top=6pt, bottom=6pt,
  title={\bfseries #1}, coltitle=white, colbacktitle=accent, fonttitle=\small
}

% \begin{jargon}{Term} definition \end{jargon}  - define a term at first use
\newtcolorbox{jargon}[1]{
  enhanced, breakable, colback=white, colframe=linegrey,
  boxrule=0.4pt, arc=2pt, left=8pt, right=8pt, top=6pt, bottom=6pt,
  title={\bfseries Jargon: #1}, coltitle=inkblue, colbacktitle=softgrey, fonttitle=\small
}

% \begin{honest}{Title} ... \end{honest}  - brutally honest finding (rule 18)
\newtcolorbox{honest}[1]{
  enhanced, breakable, colback=white, colframe=warnred,
  boxrule=0.6pt, arc=2pt, left=8pt, right=8pt, top=6pt, bottom=6pt,
  title={\bfseries #1}, coltitle=white, colbacktitle=warnred, fonttitle=\small
}

% \begin{pseudocode}{Caption} ... \end{pseudocode}
% Stands in for algorithm2e, which is NOT installed. Use \Step / \Indent inside.
\newtcolorbox{pseudocode}[1]{
  enhanced, breakable, colback=codebg, colframe=inkblue,
  boxrule=0.4pt, arc=2pt, left=10pt, right=8pt, top=6pt, bottom=6pt,
  fontupper=\ttfamily\small,
  title={\bfseries\rmfamily #1}, coltitle=white, colbacktitle=inkblue, fonttitle=\small
}
\newcommand{\Step}[1]{\par\noindent #1}
\newcommand{\Indent}[1]{\par\noindent\hspace*{2em}#1}

% ===== tikz house styles =================================================
\tikzset{
  box/.style   = {rectangle, rounded corners=2pt, draw=inkblue, fill=softgrey,
                  text=inkblue, align=center, inner sep=6pt, minimum height=9mm,
                  font=\small},
  extbox/.style= {rectangle, rounded corners=2pt, draw=linegrey, fill=white,
                  text=black!70, align=center, inner sep=6pt, minimum height=9mm,
                  font=\small, dashed},
  store/.style = {cylinder, shape border rotate=90, aspect=0.25, draw=accent,
                  fill=white, text=inkblue, align=center, inner sep=5pt,
                  font=\small},
  dec/.style   = {diamond, aspect=1.6, draw=accent, fill=white, text=inkblue,
                  align=center, inner sep=2pt, font=\scriptsize},
  flow/.style  = {-{Stealth[length=2.4mm]}, draw=inkblue, thick},
  dflow/.style = {-{Stealth[length=2.4mm]}, draw=accent, thick, dashed},
  % align=center is load-bearing: without it, a \\ inside an edge label throws
  % the misleading "perhaps a missing \item" error.
  lbl/.style   = {font=\scriptsize\itshape, text=black!65, fill=white,
                  inner sep=2pt, align=center}
}

% forest default for directory trees
\forestset{
  dirtree/.style={
    for tree={
      font=\ttfamily\small, grow'=0, child anchor=west, parent anchor=south,
      anchor=west, calign=first, inner sep=2pt,
      edge path={\noexpand\path [draw, \forestoption{edge}]
        (!u.south west) +(9pt,0) |- (.child anchor) \forestoption{edge label};},
      before typesetting nodes={if n=1{insert before={[,phantom]}}{}},
      fit=band, before computing xy={l=15pt}
    }
  }
}

% ===== misc ==============================================================
\captionsetup{font=small, labelfont={bf,color=inkblue}}
\setlist[itemize]{leftmargin=1.2em, itemsep=2pt, topsep=3pt}
\setlist[enumerate]{leftmargin=1.4em, itemsep=2pt, topsep=3pt}
\newcommand{\file}[1]{\texttt{\small #1}}
\newcommand{\code}[1]{\texttt{\small #1}}
\setcounter{tocdepth}{2}
\setcounter{secnumdepth}{3}

%   \begin{document} ... \end{document}
%
% Build:  pdflatex -interaction=nonstopmode deep-dive.tex   (run it three times, so the
%         table of contents and the TikZ positioning settle)
%
% PACKAGE NOTES. Every package used here was checked present on the machine these
% documents were built on (Debian TeX Live 2023). If you rebuild elsewhere and a build
% fails, these are the likely gaps. Do not add a package without checking first:
%   kpsewhich <pkg>.sty
%
% Deliberately NOT used, because they were absent and the build dies without them:
%   algorithm2e, algorithm, algorithmic, algpseudocode  (texlive-science not installed)
%     -> the \begin{pseudocode} environment defined below stands in for them.
%   lmodern, charter, mathptmx, times, mathpazo, courier, newtx*, tgtermes
%     -> their .sty files RESOLVE but the font files themselves are absent, so the
%        build dies with "I can't find file `bchr8t'" or similar. Verified by
%        test-compiling each one. Only Computer Modern (the default) and helvet
%        actually render. kpsewhich reporting a font it cannot render is the trap:
%        do not "upgrade" the font stack without test-compiling it first.
%   amsmath, amssymb -> not loaded, so \text{} and \blacksquare are unavailable.
%
% A fuller font stack and algorithm2e would come from:
%   sudo apt install texlive-fonts-recommended texlive-science cm-super

\usepackage[T1]{fontenc}
\usepackage[utf8]{inputenc}
% Font: Computer Modern (default). See the note above before changing this.
\usepackage[a4paper,margin=2.4cm,headheight=14pt]{geometry}
% microtype WITHOUT font expansion. This is load-bearing, do not "restore" it:
% cm-super is not installed, so T1 Computer Modern resolves to BITMAP fonts, and
% pdfTeX's auto-expansion needs scalable ones. With expansion on, any document
% long enough to reach a second page dies with
%   "pdfTeX error (font expansion): auto expansion is only possible with scalable fonts"
% Protrusion (the bigger visual win) still works and stays on.
% `sudo apt install cm-super` would lift this; until then, leave expansion off.
\usepackage[protrusion=true,expansion=false]{microtype}
% Kill hyphen ligatures in TYPEWRITER text. This is load-bearing, not cosmetic. Under T1,
% \texttt{--verbose} silently renders as "-verbose", and \texttt{---} as "--", so a
% document can print a command-line flag as the WRONG flag and look perfectly fine doing
% it. Verified by rendering the page and reading the glyphs. lstlisting and \verb were
% always safe; this makes \code, \file and \texttt safe too, which is what prose uses.
% Ordinary prose is unaffected and still ligates, which is what you want there.
\DisableLigatures[-]{family=tt*}
\usepackage{graphicx}
\usepackage{xcolor}
\usepackage{booktabs}
\usepackage{enumitem}
\usepackage{fancyhdr}
\usepackage{titlesec}
\usepackage{tcolorbox}
\usepackage{listings}
\usepackage{float}
\usepackage{caption}
\usepackage{tikz}
\usepackage{forest}
\usepackage{pgfplots}
\pgfplotsset{compat=1.18}
\usetikzlibrary{arrows.meta,positioning,shapes.geometric,fit,backgrounds,calc,decorations.pathreplacing}
\tcbuselibrary{skins,breakable}
\usepackage[hidelinks,bookmarks=true,bookmarksnumbered=true]{hyperref}

% ===== palette ===========================================================
\definecolor{inkblue}{HTML}{1F3A5F}
\definecolor{accent}{HTML}{2E7DAF}
\definecolor{softgrey}{HTML}{F4F6F8}
\definecolor{linegrey}{HTML}{C8D0D8}
\definecolor{codebg}{HTML}{F7F7F4}
\definecolor{warnred}{HTML}{A33B3B}
\definecolor{okgreen}{HTML}{2E7D53}

% ===== headings ==========================================================
\titleformat{\section}{\Large\bfseries\color{inkblue}}{\thesection}{0.7em}{}
\titleformat{\subsection}{\large\bfseries\color{inkblue}}{\thesubsection}{0.6em}{}
\titleformat{\subsubsection}{\normalsize\bfseries\color{accent}}{\thesubsubsection}{0.5em}{}

\pagestyle{fancy}
\fancyhf{}
\fancyhead[L]{\small\color{inkblue}\nouppercase{\leftmark}}
\fancyhead[R]{\small\color{inkblue}\thepage}
\renewcommand{\headrulewidth}{0.4pt}
\renewcommand{\headrule}{\hbox to\headwidth{\color{linegrey}\leaders\hrule height \headrulewidth\hfill}}

% ===== code ==============================================================
\lstdefinestyle{house}{
  backgroundcolor=\color{codebg},
  basicstyle=\ttfamily\footnotesize,
  keywordstyle=\color{inkblue}\bfseries,
  commentstyle=\color{gray}\itshape,
  stringstyle=\color{okgreen},
  numbers=left, numberstyle=\tiny\color{gray}, numbersep=8pt,
  frame=single, rulecolor=\color{linegrey},
  breaklines=true, breakatwhitespace=false,
  showstringspaces=false, tabsize=2,
  captionpos=b, xleftmargin=14pt, framexleftmargin=14pt,
  columns=fullflexible, keepspaces=true,
  upquote=true,  % render ' and " as straight quotes, not curly
  % NON-ASCII INSIDE A LISTING IS A HARD BUILD FAILURE, not a cosmetic issue: listings
  % reads raw bytes, so inputenc dies on the first multi-byte character ("Invalid UTF-8
  % byte sequence") and the whole document fails to build. These mappings cover what
  % realistically turns up in source code: Latin-1 accents, Spanish punctuation, smart
  % quotes picked up from a word processor, the degree and section signs.
  % NOT COVERED, and it cannot be here: Urdu, Arabic and Balochi script. pdflatex has no
  % font for them. Describe or transliterate such text; never paste it into a listing.
  literate=%
    {á}{{\'a}}1 {é}{{\'e}}1 {í}{{\'i}}1 {ó}{{\'o}}1 {ú}{{\'u}}1
    {Á}{{\'A}}1 {É}{{\'E}}1 {Í}{{\'I}}1 {Ó}{{\'O}}1 {Ú}{{\'U}}1
    {à}{{\`a}}1 {è}{{\`e}}1 {ì}{{\`i}}1 {ò}{{\`o}}1 {ù}{{\`u}}1
    {ä}{{\"a}}1 {ë}{{\"e}}1 {ï}{{\"i}}1 {ö}{{\"o}}1 {ü}{{\"u}}1
    {Ä}{{\"A}}1 {Ö}{{\"O}}1 {Ü}{{\"U}}1 {ß}{{\ss}}1
    {ñ}{{\~n}}1 {Ñ}{{\~N}}1 {ç}{{\c c}}1 {Ç}{{\c C}}1
    {¿}{{\textquestiondown}}1 {¡}{{\textexclamdown}}1
    {°}{{\textdegree}}1 {§}{{\S}}1 {£}{{\pounds}}1 {€}{{EUR}}1
    {“}{{\textquotedblleft}}1 {”}{{\textquotedblright}}1
    {‘}{{\textquoteleft}}1 {’}{{\textquoteright}}1
    {–}{{-}}1 {—}{{-}}1 {…}{{...}}1 {→}{{$\rightarrow$}}1 {✓}{{x}}1
}
\lstset{style=house}

% ===== callout boxes =====================================================
% \begin{keyidea}{Title} ... \end{keyidea}   - a concept a beginner must grasp
\newtcolorbox{keyidea}[1]{
  enhanced, breakable, colback=softgrey, colframe=accent,
  boxrule=0.4pt, arc=2pt, left=8pt, right=8pt, top=6pt, bottom=6pt,
  title={\bfseries #1}, coltitle=white, colbacktitle=accent, fonttitle=\small
}

% \begin{jargon}{Term} definition \end{jargon}  - define a term at first use
\newtcolorbox{jargon}[1]{
  enhanced, breakable, colback=white, colframe=linegrey,
  boxrule=0.4pt, arc=2pt, left=8pt, right=8pt, top=6pt, bottom=6pt,
  title={\bfseries Jargon: #1}, coltitle=inkblue, colbacktitle=softgrey, fonttitle=\small
}

% \begin{honest}{Title} ... \end{honest}  - brutally honest finding (rule 18)
\newtcolorbox{honest}[1]{
  enhanced, breakable, colback=white, colframe=warnred,
  boxrule=0.6pt, arc=2pt, left=8pt, right=8pt, top=6pt, bottom=6pt,
  title={\bfseries #1}, coltitle=white, colbacktitle=warnred, fonttitle=\small
}

% \begin{pseudocode}{Caption} ... \end{pseudocode}
% Stands in for algorithm2e, which is NOT installed. Use \Step / \Indent inside.
\newtcolorbox{pseudocode}[1]{
  enhanced, breakable, colback=codebg, colframe=inkblue,
  boxrule=0.4pt, arc=2pt, left=10pt, right=8pt, top=6pt, bottom=6pt,
  fontupper=\ttfamily\small,
  title={\bfseries\rmfamily #1}, coltitle=white, colbacktitle=inkblue, fonttitle=\small
}
\newcommand{\Step}[1]{\par\noindent #1}
\newcommand{\Indent}[1]{\par\noindent\hspace*{2em}#1}

% ===== tikz house styles =================================================
\tikzset{
  box/.style   = {rectangle, rounded corners=2pt, draw=inkblue, fill=softgrey,
                  text=inkblue, align=center, inner sep=6pt, minimum height=9mm,
                  font=\small},
  extbox/.style= {rectangle, rounded corners=2pt, draw=linegrey, fill=white,
                  text=black!70, align=center, inner sep=6pt, minimum height=9mm,
                  font=\small, dashed},
  store/.style = {cylinder, shape border rotate=90, aspect=0.25, draw=accent,
                  fill=white, text=inkblue, align=center, inner sep=5pt,
                  font=\small},
  dec/.style   = {diamond, aspect=1.6, draw=accent, fill=white, text=inkblue,
                  align=center, inner sep=2pt, font=\scriptsize},
  flow/.style  = {-{Stealth[length=2.4mm]}, draw=inkblue, thick},
  dflow/.style = {-{Stealth[length=2.4mm]}, draw=accent, thick, dashed},
  % align=center is load-bearing: without it, a \\ inside an edge label throws
  % the misleading "perhaps a missing \item" error.
  lbl/.style   = {font=\scriptsize\itshape, text=black!65, fill=white,
                  inner sep=2pt, align=center}
}

% forest default for directory trees
\forestset{
  dirtree/.style={
    for tree={
      font=\ttfamily\small, grow'=0, child anchor=west, parent anchor=south,
      anchor=west, calign=first, inner sep=2pt,
      edge path={\noexpand\path [draw, \forestoption{edge}]
        (!u.south west) +(9pt,0) |- (.child anchor) \forestoption{edge label};},
      before typesetting nodes={if n=1{insert before={[,phantom]}}{}},
      fit=band, before computing xy={l=15pt}
    }
  }
}

% ===== misc ==============================================================
\captionsetup{font=small, labelfont={bf,color=inkblue}}
\setlist[itemize]{leftmargin=1.2em, itemsep=2pt, topsep=3pt}
\setlist[enumerate]{leftmargin=1.4em, itemsep=2pt, topsep=3pt}
\newcommand{\file}[1]{\texttt{\small #1}}
\newcommand{\code}[1]{\texttt{\small #1}}
\setcounter{tocdepth}{2}
\setcounter{secnumdepth}{3}


\title{\textbf{Trade Intelligence Copilot}\\[4pt]
\large A Deep Dive: agentic retrieval over trade data,\\with an evaluation layer that is the point}
\author{Explainer for Salman Adnan's portfolio}
\date{}

\begin{document}
\maketitle
\thispagestyle{empty}

\begin{keyidea}{Read this first: what this document assumes}
Nothing. This document assumes you have never built a retrieval system, never
written a line of SQL, and have never heard the phrase ``agentic retrieval''.
Every technical term is defined in a \textbf{Jargon} box the first time it
appears. If you already know the vocabulary, skip the boxes and read the
diagrams.
\end{keyidea}

\tableofcontents
\newpage

% =====================================================================
\section{What this project is, in one paragraph}

Trade Intelligence Copilot answers plain-English questions about Pakistani
sugar-sector import data and about trade-regulation concepts. You type
``Which country was the largest source of imports by value?'' and it replies
with a number taken from a real database, showing you the exact query it ran
and the rows it got back. You type ``What is a phytosanitary certificate?''
and it replies with a definition drawn from a specific document, and tells you
which document. You type ``What were Pakistan's wheat exports to India in
2019?'' and it says: \emph{this question is not supported by the available
data}, because the dataset holds no exports, no wheat, and no 2019.

That last behaviour is the hard one, and it is why the project exists.

\begin{keyidea}{The thesis of the project}
Most demonstrations of this technology answer confidently and are never
measured. This one ships a 28-item test set with ground-truth answers computed
directly from the database, runs the whole system against a live model API,
publishes the real numbers, and then checks its own automatic grader against
independent labels (11 of 12 agree). The headline is not the copilot. The headline is the
\textbf{evaluation layer}.
\end{keyidea}

% =====================================================================
\section{The ideas, from first principles}

Before any code, here is the vocabulary. Each idea is small; the system is just
these ideas wired together.

\subsection{Large language models, and their central weakness}

\begin{jargon}{Large language model (LLM)}
A program trained on an enormous amount of text to predict what text should
come next. Given ``The capital of France is'', it answers ``Paris''. Modern
ones follow instructions, write code, and produce structured output. This
project uses Google's Gemini models through Vertex AI, Google's cloud platform
for machine-learning services.
\end{jargon}

An LLM knows only what it absorbed during training. It does not know what is in
\emph{your} database. Worse, when it does not know, it often invents an answer
that reads perfectly. This is called \textbf{hallucination}, and for a system
that reports trade values in dollars it is fatal: a confidently invented
\$175 million is worse than no answer at all.

\begin{jargon}{Hallucination}
When a language model states something fluent, plausible and false. The model is
not lying; it has no concept of truth. It is producing text that pattern-matches
what a correct answer looks like.
\end{jargon}

\subsection{Retrieval-augmented generation (RAG)}

The standard fix is to stop asking the model to \emph{remember} and start
forcing it to \emph{read}.

\begin{jargon}{Retrieval-augmented generation (RAG)}
A pattern in two beats. First \textbf{retrieve}: go and fetch the specific facts
or documents relevant to the question from a trusted source. Then
\textbf{generate}: hand the model \emph{only} those facts and instruct it to
answer using nothing else. The model contributes fluency and reasoning; the
trusted source contributes the facts.
\end{jargon}

This is the difference between a closed-book exam and an open-book exam where
you are only allowed to cite the page in front of you.

\subsection{Agentic retrieval: the step that makes it ``agentic''}

Plain RAG always does the same thing: search a pile of documents, paste the
results in, answer. But the questions here are not all the same shape. ``What
was the total USD value of all imports?'' is not a question about documents at
all; it is arithmetic over a database. ``What is a certificate of origin?'' is
purely a document question. ``How much sugar was imported, and why is sugar
policy-sensitive?'' is both at once.

\begin{jargon}{Agentic retrieval}
Retrieval where the system first \textbf{decides for itself} which sources to
consult, and how many, rather than always running one fixed search. A model
reads the question and picks the tools. That decision step is what earns the
word ``agentic''. The system is choosing a strategy, not executing a fixed
recipe.
\end{jargon}

\begin{jargon}{Tool}
A concrete capability the system can invoke: run a database query, search
documents, look up a code. The model does not perform the tool's work; it
\emph{names} the tool it wants and ordinary code performs the work. This
distinction matters for safety: the model never touches the database directly.
\end{jargon}

In this project the planner is an LLM whose entire job is to answer: does this
question need SQL, document retrieval, an HS-code lookup, or some combination?
It is explicitly told \emph{not} to answer the question itself.

\subsection{Text-to-SQL}

\begin{jargon}{SQL}
Structured Query Language, the standard language for asking questions of a
relational database (data in tables of rows and columns).
\code{SELECT COUNT(*) FROM trade\_transactions} means ``count every row in the
trade\_transactions table''. \code{SELECT} reads data; other commands such as
\code{INSERT}, \code{UPDATE}, \code{DELETE} and \code{DROP} change or destroy
it.
\end{jargon}

\begin{jargon}{Text-to-SQL}
Having a language model translate a plain-English question into a SQL query,
then running that query on the real database and using the returned rows as the
answer. The model supplies the translation; the database supplies the number.
\end{jargon}

This is the single most important idea for trustworthiness here. When you ask
for a total, the model is never asked to recall or estimate the total. It is
asked to write a query. The database computes the number. If the query is
right, the number is right by construction.

The obvious danger: the model is writing executable code against your database,
and a model can be talked into writing \code{DROP TABLE}. Section~\ref{sec:guard}
is entirely about containing that.

\subsection{Embeddings, vectors and cosine similarity}

For the document half of the system, we need to find text that is
\emph{about} the same thing as the question, even when it shares no words with
it. Searching for the literal word ``phytosanitary'' is brittle; the user might
ask about ``plant health paperwork''.

\begin{jargon}{Embedding}
A list of numbers (here, 768 of them) that represents the \emph{meaning} of a
piece of text, produced by a model trained for the purpose. The useful property:
texts with similar meanings get similar lists, even with no words in common.
\end{jargon}

\begin{jargon}{Vector}
A list of numbers treated as a point, or an arrow from the origin, in space.
768 numbers means a point in 768-dimensional space. You cannot picture it; the
arithmetic works exactly as it does in the two dimensions you can picture.
\end{jargon}

\begin{jargon}{Cosine similarity}
A measure of whether two vectors point the same way, ignoring their length.
It runs from $1.0$ (identical direction, so: same meaning) through $0$
(unrelated) to $-1.0$ (opposite). ``Similar meaning'' becomes ``small angle
between the arrows'', and that is a number you can sort by.
\end{jargon}

The trick that makes this cheap: if every vector is first scaled to length
exactly $1$ (``normalized''), then cosine similarity is just the dot product,
which is one multiply-and-add per dimension. The project exploits this
directly, and Section~\ref{sec:vectorstore} shows the four lines of code where
it happens.

\subsection{Chunks and citation}

\begin{jargon}{Chunk}
A document cut into a retrievable piece. You do not embed a whole document,
because then a hit tells the reader ``the answer is somewhere in these five
pages''. This project splits each document at its \code{\#\#} section headings,
so one chunk is one whole section, which is also the natural thing to cite.
\end{jargon}

Each chunk carries an id like \code{phytosanitary-sps-basics\#s2}, meaning
section 2 of that document. When the copilot states a rule, it prints that id.
The claim is auditable: you can go read the exact section it came from.

\subsection{An evaluation layer, and why it is the whole point}

\begin{jargon}{Evaluation layer (``eval'')}
Machinery that measures whether the system is actually right, on a fixed set of
questions with known correct answers, producing numbers you can publish and
argue with. Not a demo, not a vibe check: a score.
\end{jargon}

\begin{jargon}{Ground truth}
The known-correct answer, established independently of the system being tested.
Here, the gold answer to ``how many shipments came from China?'' is not a human
guess: it is the result of running the counting query against the database
directly. It cannot be wrong, so ``correct'' stops being a matter of opinion.
\end{jargon}

\begin{keyidea}{Why an eval changes the argument}
Without an eval you can say ``it seems to work''. An interviewer cannot check
that, and neither can you. With an eval you can say ``on 15 questions whose
answers I computed from the database, it got 15 right, and here is the
harness''. It also catches bugs no amount of looking at pretty answers ever
would: this project's text-to-SQL step once produced a query that returned zero
rows, and the answer that reported ``0 shipments'' looked completely
reasonable. Only the gold number, which was dozens, exposed it.
\end{keyidea}

\begin{jargon}{LLM-as-judge}
Using a second, usually stronger, language model to grade the first one's
output. Here \code{gemini-2.5-pro} reads the answer plus the evidence and rules
on whether every claim is supported. Necessary because ``did this paragraph
invent anything?'' cannot be checked with string matching.
\end{jargon}

\begin{jargon}{Faithfulness}
Whether an answer is fully supported by the evidence it was given. Note
carefully that faithfulness is \textbf{not} correctness. An answer that
faithfully reports the output of a wrong query is faithful and wrong. This
project keeps the two measurements separate on purpose, and
Section~\ref{sec:findings} explains why that separation was forced on it by a
real bug.
\end{jargon}

\begin{keyidea}{A judge you have not validated is not evidence}
If you grade a model with a model, who grades the judge? This project answers
that: a 12-item subset is labelled independently, each answer read against its
evidence, and the harness reports how often the judge agrees with those labels (11
of 12, on both the first version and the current run). That agreement number is
what makes the faithfulness number mean anything at all.
\end{keyidea}

% =====================================================================
\section{Map of the repository}

Every tracked file, and what it is for.

\begin{figure}[H]
\centering
\begin{forest} dirtree
[trade-intelligence-copilot
  [app/ {\rmfamily\itshape the copilot itself}
    [config.py {\rmfamily\itshape all settings + secrets read from environment}]
    [llm.py {\rmfamily\itshape Gemini wrapper + token-cost meter}]
    [embeddings.py {\rmfamily\itshape Vertex embedding REST calls}]
    [vectorstore.py {\rmfamily\itshape numpy cosine-similarity search}]
    [schema\_prompt.py {\rmfamily\itshape DB schema + scope facts + SQL rules}]
    [planner.py {\rmfamily\itshape THE AGENT: picks the tools}]
    [tools.py {\rmfamily\itshape text-to-SQL, retrieval, hs\_lookup}]
    [db.py {\rmfamily\itshape SQLite access + the SQL safety guard}]
    [answer.py {\rmfamily\itshape grounded synthesis + refusal}]
    [agent.py {\rmfamily\itshape orchestration, trace, cost}]
    [server.py {\rmfamily\itshape FastAPI: /ask, /health, /}]
    [templates/index.html {\rmfamily\itshape single-page UI}]
  ]
  [data/
    [build\_db.py {\rmfamily\itshape seeded generator -> trade.sqlite}]
    [trade.sqlite {\rmfamily\itshape 2,000 synthetic shipments (prebuilt, shipped)}]
    [regulations/ {\rmfamily\itshape 16 authored explainer documents}]
    [chunks.json {\rmfamily\itshape 48 chunks + metadata (prebuilt)}]
    [vectors.npz {\rmfamily\itshape 48x768 normalized matrix (prebuilt)}]
  ]
  [ingest/ingest.py {\rmfamily\itshape chunk + embed the corpus}]
  [eval/ {\rmfamily\itshape THE HEADLINE}
    [testset.jsonl {\rmfamily\itshape 28 questions + gold answers}]
    [run\_eval.py {\rmfamily\itshape harness: run, judge, score}]
    [transcripts.json {\rmfamily\itshape every answer from the real run}]
    [results.json {\rmfamily\itshape the published metrics}]
    [human\_labels.json {\rmfamily\itshape 12 independent labels for the judge check}]
  ]
  [tests/ {\rmfamily\itshape 33 tests: 29 offline, 4 live}]
  [docs/ {\rmfamily\itshape ARCHITECTURE.md, EVALUATION.md}]
  [requirements.txt]
  [.env.example {\rmfamily\itshape config template, no real values}]
]
\end{forest}
\caption{The repository. Roughly 1,000 lines of Python, deliberately small: the
substance is in \file{eval/}, not in framework scaffolding.}
\end{figure}

\begin{keyidea}{No framework}
There is no LangChain, no LlamaIndex, no agent framework. The planner, the tool
dispatch, the vector store and the eval harness are all hand-written and total
about a thousand lines. This is a defensible choice for a project whose purpose
is to demonstrate that the author understands the mechanism: a framework would
hide exactly the parts worth showing.
\end{keyidea}

% =====================================================================
\section{End-to-end architecture}

Here is the entire system. One question enters at the left; an answer with
citations, a trace, a cost and a latency leaves at the right.

\begin{figure}[H]
\centering
\resizebox{\textwidth}{!}{
\begin{tikzpicture}[node distance=8mm]
  \node[box] (q) {Question\\\scriptsize ``How many shipments\\\scriptsize came from China?''};
  \node[box, right=14mm of q, fill=accent!12] (plan) {\textbf{Planner}\\\scriptsize gemini-2.5-flash\\\scriptsize picks the tools};

  \node[box, above right=6mm and 16mm of plan] (sql) {\textbf{text-to-SQL}\\\scriptsize LLM writes SELECT};
  \node[box, right=16mm of plan] (ret) {\textbf{retrieval}\\\scriptsize embed + cosine};
  \node[box, below right=6mm and 16mm of plan] (hs) {\textbf{hs\_lookup}\\\scriptsize direct SQL join};

  \node[box, right=12mm of sql, fill=warnred!12] (guard) {\textbf{SQL guard}\\\scriptsize single SELECT only};
  \node[store, right=12mm of guard] (db) {SQLite\\\scriptsize read-only};
  \node[store, right=12mm of ret] (vs) {vectors\\\scriptsize 48x768};
  \node[store, right=12mm of hs] (db2) {SQLite\\\scriptsize hierarchy};

  \node[dec, right=44mm of ret, xshift=14mm] (ev) {any\\evidence?};
  \node[box, above right=8mm and 12mm of ev, fill=warnred!12] (refuse) {\textbf{Refuse}\\\scriptsize deterministic,\\\scriptsize costs nothing};
  \node[box, below right=8mm and 12mm of ev, fill=accent!12] (syn) {\textbf{Synthesizer}\\\scriptsize gemini-2.5-flash\\\scriptsize evidence only};
  \node[box, right=14mm of ev, xshift=30mm] (out) {\{answer, supported,\\citations, trace,\\cost, latency\}};

  \draw[flow] (q) -- (plan);
  \draw[flow] (plan) -- (sql); \draw[flow] (plan) -- (ret); \draw[flow] (plan) -- (hs);
  \draw[flow] (sql) -- (guard); \draw[flow] (guard) -- (db);
  \draw[flow] (ret) -- (vs); \draw[flow] (hs) -- (db2);
  \draw[flow] (db) -- (ev); \draw[flow] (vs) -- (ev); \draw[flow] (db2) -- (ev);
  \draw[flow] (ev) -- node[lbl,above]{no} (refuse);
  \draw[flow] (ev) -- node[lbl,below]{yes} (syn);
  \draw[flow] (refuse) -- (out); \draw[flow] (syn) -- (out);
\end{tikzpicture}}
\caption{The agent loop. Three model calls on a typical answered question
(plan, text-to-SQL, synthesize); two on a refusal, because the deterministic
check fires before the expensive synthesis call.}
\label{fig:arch}
\end{figure}

Read that diagram as a sentence: \emph{an LLM decides which tools to run;
ordinary code runs them; a cheap check refuses when they returned nothing; and
a second LLM writes prose using only what they returned.}

\begin{keyidea}{Where the trust lives}
Notice what the models are and are not allowed to do. The planner chooses but
does not answer. The text-to-SQL model writes a query but never executes it: the
guard and the read-only connection do that. The synthesizer writes prose but is
handed a fixed block of evidence and told to use nothing else. No model is ever
trusted with both the facts and the final word.
\end{keyidea}

% =====================================================================
\section{The data layer}

\subsection{Where the trade data comes from}

\file{data/build\_db.py} \emph{generates} the trade data: a compact SQLite
database of synthetic Pakistani sugar-sector imports. Every importer, supplier,
shipping agent, price and quantity is made up. Only the public scaffolding is
real: the HS chapter 17 headings and the country names. The first version of the
project ran on a dataset that cannot be published, and this generator replaced it.

\begin{jargon}{SQLite}
A complete relational database that is just a single file on disk, with no
server to install or run. Ideal here: the demo database ships inside the repo
and the project runs with nothing to set up.
\end{jargon}

\begin{jargon}{Seeded generator}
A program that draws ``random'' values from a pseudo-random number generator
started at a fixed number, the seed. Same seed, same sequence, so the output is
identical on every run. Here \code{SEED = 1717}, and two builds of
\file{trade.sqlite} are byte-for-byte the same file.
\end{jargon}

The script is deliberately plain: drop the file if it exists, create the schema,
write the HS hierarchy, then draw 2,000 shipments. Each shipment picks a
subcategory by weight, an origin country (bulk sugar and molasses only from
Brazil, Thailand or India, weighted so Brazil clearly leads by value), a price
from that subcategory's band, and a log-uniform quantity, so there are many small
lots and a few large ones. Determinism is what makes the evaluation honest after
the switch: \file{eval/build\_testset.py} recomputes every gold answer from the
generated file with a query, so nobody types a gold number by hand.

\subsection{The schema}

\begin{figure}[H]
\centering
\resizebox{0.95\textwidth}{!}{
\begin{tikzpicture}[node distance=9mm]
  \node[box] (p) {\textbf{products}\\\scriptsize id, hs\_code, name\\\scriptsize e.g. 1701 `Cane or beet sugar'};
  \node[box, below=of p] (c) {\textbf{product\_categories}\\\scriptsize id, hs\_code, name, product\_id};
  \node[box, below=of c] (s) {\textbf{product\_subcategories}\\\scriptsize id, hs\_code, name, category\_id};
  \node[box, below=of s] (i) {\textbf{product\_items}\\\scriptsize id, name, subcategory\_id};
  \node[box, right=42mm of s, fill=accent!12] (t) {\textbf{trade\_transactions}\\\scriptsize 2,000 rows, ALL imports\\\scriptsize country, hs\_code, buyer, seller\\\scriptsize usd, pkr, qty\_kg, qty\_mt\\\scriptsize reporting\_date 2025-01-01 .. 2025-12-31};
  \node[box, above=14mm of t] (co) {\textbf{companies}\\\scriptsize 28 rows, a directory\\\scriptsize not fully linked to transactions};

  \draw[flow] (c) -- node[lbl,right]{product\_id} (p);
  \draw[flow] (s) -- node[lbl,right]{category\_id} (c);
  \draw[flow] (i) -- node[lbl,right]{subcategory\_id} (s);
  \draw[flow] (t) -- node[lbl,above]{product\_item\_id} (i);
\end{tikzpicture}}
\caption{The HS-code hierarchy on the left, the fact table on the right. Four
indexes exist on \code{trade\_transactions}: country, hs\_code, trade\_type and
reporting\_date, which are exactly the columns the questions filter on.}
\label{fig:schema}
\end{figure}

\begin{jargon}{HS code}
The Harmonized System, an international product-classification standard. The
first two digits are the chapter, the first four the heading. This dataset is
entirely chapter 17: 1701 cane and beet sugar, 1702 other sugars and syrups,
1703 molasses, 1704 sugar confectionery. Codes are stored as \emph{text}, not
numbers, so that \code{1702.1100} keeps its shape and prefix matching works.
\end{jargon}

\subsection{The scope facts: the quiet foundation of honest refusal}

\file{app/schema\_prompt.py} carries three blocks of text that are injected into
prompts. The second, \code{DATASET\_SCOPE}, is the most interesting thing in the
file, and it is easy to overlook:

\begin{lstlisting}[language=Python,caption={The scope facts, given to both the planner and the text-to-SQL step (abridged).}]
DATASET_SCOPE = """
DATASET SCOPE (important for deciding if a question is answerable):
- trade_transactions contains 2000 rows. Every row is an IMPORT.
  There are NO export rows in this dataset.
- All shipments are in HS chapter 17 (sugar and sugar products).
  There is no data for other commodities (no rice, wheat, cotton...).
- reporting_date ranges from 2025-01-01 to 2025-12-31 only.
...
If a question asks about data outside this scope (exports, non-sugar
commodities, other years...), it CANNOT be answered from the data.
"""
\end{lstlisting}

\begin{keyidea}{Refusal starts here, not in the refusal code}
A model cannot decline a question about exports unless it knows the dataset has
no exports. These sentences are what let the planner recognise
``Pakistan's wheat exports in 2019'' as three separate impossibilities at once.
Every one of these facts is checkable against the database, and I checked all of
them (Section~\ref{sec:verify}): 2,000 rows, \code{trade\_type} is
\code{'import'} for every row, dates run exactly 2025-01-01 to 2025-12-31,
28 companies. The docstring calls them ``computed directly from the database'',
and that is true.
\end{keyidea}

The third block, \code{TEXT\_TO\_SQL\_RULES}, is where a real bug was fixed. It
now insists on wildcard matching for any name or free-text column, and forbids
matching a name with \code{=}:

\begin{lstlisting}[caption={One rule from TEXT\_TO\_SQL\_RULES, added after the bug in Section 13. The example name is deliberately not one the evaluation asks about.}]
- Match names or free text with LIKE wrapped in % on both sides, for example
  buyer LIKE '%Tamrin%' or country LIKE '%South Africa%'. Never match a name with =.
\end{lstlisting}

\subsection{The regulations corpus}

Sixteen short Markdown documents in \file{data/regulations/}, covering HS-code
basics, phytosanitary and SPS rules, customs-duty concepts, import
documentation, anti-dumping and safeguards, letters of credit, the Pakistan
Single Window, and one document per HS heading 1701 to 1704.

\begin{honest}{These are authored summaries, not law, and the project says so}
The corpus is plain-language explainers written for this demo. They reproduce no
statute, no SRO number and no actual duty rate, and \file{\_README.md} in that
folder states this plainly, as does \file{docs/EVALUATION.md}. This is an
honest and necessary disclosure: the faithfulness metric measures grounding in
\emph{this corpus}, not legal accuracy. A retrieval system that is 100\%
faithful to a made-up document is 100\% useless for actual customs compliance.
The project does not overclaim here, and it should not start.
\end{honest}

\subsection{Ingestion: chunking and embedding}

\begin{sloppypar}
\file{ingest/ingest.py} turns those documents into searchable vectors. Its
cleverest function, \code{split\_sections}, walks a document line by line. It
treats every second-level Markdown heading (a line starting with two hash
characters) as a chunk boundary, and captures the top-level heading as the
document title.
\end{sloppypar}

\begin{pseudocode}{ingest.run(), the whole pipeline}
\Step{1. refuse to start unless credentials are present}
\Step{2. for each *.md in regulations/, skipping files starting with "\_":}
\Indent{read it; find the "\# " document title}
\Indent{split on "\#\# " headings into (section\_title, section\_text)}
\Indent{skip the front-matter block (section title == doc title)}
\Indent{chunk id = "<filename stem>\#s<index>"}
\Indent{embed\_text = "<doc title> :: <section title>\textbackslash n<section text>"}
\Step{3. embed every embed\_text in batches of 5 (RETRIEVAL\_DOCUMENT)}
\Step{4. normalize all vectors to unit length}
\Step{5. write vectors.npz (matrix) + chunks.json (metadata, embed\_text dropped)}
\end{pseudocode}

Two details worth pausing on, because both are real craft:

\begin{itemize}
\item \textbf{The embedded text is not the stored text.} What gets embedded is
  \code{"<doc title> :: <section title>\textbackslash n<section text>"}, so the
  document and section headings contribute to the meaning of the vector. A
  section called ``2. Reading the digits'' is nearly meaningless on its own; as
  ``HS codes: structure and how classification works :: 2. Reading the digits''
  it is findable. What is \emph{stored} and later shown to the reader is the
  plain section text.
\item \textbf{The task type is set.} Vertex embedding models accept a
  \code{task\_type}. Ingestion passes \code{RETRIEVAL\_DOCUMENT}; a live query
  passes \code{RETRIEVAL\_QUERY} (\file{app/tools.py}, line 51). The model
  embeds a question and a passage into the same space but with the asymmetry
  between them accounted for. Getting this backwards silently degrades every
  search, and it is a common mistake.
\end{itemize}

The result, verified: 48 chunks from 16 documents, each a 768-dimensional
vector, cached on disk so that neither the API nor the eval ever re-embeds the
corpus.

% =====================================================================
\section{Embeddings and the vector store}
\label{sec:vectorstore}

\subsection{The retrieval path}

\begin{figure}[H]
\centering
\resizebox{\textwidth}{!}{
\begin{tikzpicture}[node distance=7mm]
  \node[box] (q) {retrieval\_query\\\scriptsize ``what is a phytosanitary\\\scriptsize certificate''};
  \node[box, right=12mm of q, fill=accent!12] (emb) {\textbf{embed\_one()}\\\scriptsize Vertex predict REST\\\scriptsize task=RETRIEVAL\_QUERY};
  \node[box, right=12mm of emb] (qv) {query vector\\\scriptsize 768 floats};
  \node[box, right=12mm of qv] (norm) {normalize\\\scriptsize length := 1};
  \node[store, above right=8mm and 12mm of norm] (mat) {matrix\\\scriptsize 48 x 768\\\scriptsize pre-normalized};
  \node[box, below right=8mm and 12mm of norm, fill=accent!12] (dot) {\textbf{matrix @ q}\\\scriptsize 48 dot products\\\scriptsize = cosine scores};
  \node[box, right=14mm of dot] (sort) {argsort\\\scriptsize take top 4};
  \node[dec, right=12mm of sort] (thr) {score\\$\geq 0.55$?};
  \node[box, above right=6mm and 12mm of thr] (keep) {\textbf{strong}\\\scriptsize -> evidence};
  \node[box, below right=6mm and 12mm of thr, fill=warnred!12] (drop) {\textbf{dropped}\\\scriptsize too weak to cite};

  \draw[flow] (q) -- (emb); \draw[flow] (emb) -- (qv); \draw[flow] (qv) -- (norm);
  \draw[dflow] (mat) -- (dot);
  \draw[flow] (norm) -- (dot);
  \draw[flow] (dot) -- (sort); \draw[flow] (sort) -- (thr);
  \draw[flow] (thr) -- node[lbl,above]{yes} (keep);
  \draw[flow] (thr) -- node[lbl,below]{no} (drop);
\end{tikzpicture}}
\caption{The retrieval path. The threshold at the end is what makes refusal
possible: retrieval always returns its top 4, so ``nothing relevant'' has to be
\emph{defined}, and it is defined as ``nothing scored 0.55 or better''.}
\label{fig:retrieval}
\end{figure}

\begin{keyidea}{Why the threshold is load-bearing}
A nearest-neighbour search never fails. Ask it about wheat exports and it will
cheerfully hand back the four \emph{least irrelevant} sugar documents, with
scores like 0.31. Without a threshold those become ``evidence'' and the
synthesizer is invited to hallucinate on top of them. \code{RETRIEVAL\_MIN\_SCORE
= 0.55} converts a ranking into a judgment. The README records that an earlier
version answered out-of-scope questions precisely because weak chunks were
allowed through.
\end{keyidea}

\subsection{The store itself}

\file{app/vectorstore.py} is 57 lines and needs no more. The search is four
lines of arithmetic:

\begin{lstlisting}[language=Python,caption={The entire search, from app/vectorstore.py.}]
q = np.asarray(query_vec, dtype=np.float32)
n = np.linalg.norm(q)
if n:
    q = q / n            # normalize the query
scores = self.matrix @ q # 48 cosine similarities, one matrix-vector product
order = np.argsort(-scores)[:top_k]
\end{lstlisting}

Because the stored matrix rows were normalized at save time and the query is
normalized here, \code{matrix @ q} \emph{is} the vector of cosine similarities.
The negation inside \code{argsort} is the standard trick for sorting descending.
This is an \textbf{exact} search: it compares against all 48 chunks and cannot
miss the true best match.

\begin{keyidea}{The pgvector decision, defended}
The obvious choice was pgvector, a PostgreSQL extension for vector search. The
target machine has PostgreSQL 18 without the \code{vector} extension, no Docker,
and no superuser path to build it. Rather than fight the infrastructure, the
author wrote the 57 lines above. This is the correct call and should be argued
confidently: for 48 chunks, brute force is instant and \emph{exact}, while an
approximate-nearest-neighbour index would be slower to build, approximate, and a
whole extra moving part. The honest trade-off, stated in
\file{docs/ARCHITECTURE.md}: no ANN index and no server-side filtering, which
would begin to matter at perhaps 100,000 chunks, not 48.
\end{keyidea}

\subsection{Why embeddings go through raw REST}

\file{app/embeddings.py} does not use the \code{google-genai} SDK that the rest
of the project uses for generation. Its docstring explains: the SDK does not
expose \code{text-multilingual-embedding-002} uniformly in Vertex mode, so the
module calls the \code{:predict} endpoint directly, minting a bearer token from
the service-account credentials via \code{google.auth}.

\begin{jargon}{Service account and bearer token}
A service account is a non-human identity for a program, with a private key in a
JSON file. The code exchanges that key for a short-lived \textbf{bearer token},
a string sent in the \code{Authorization} header that proves the request is
authorised. Tokens expire, hence the \code{if not \_creds.valid: refresh(...)}.
\end{jargon}

\begin{keyidea}{Nothing sensitive is in this repository}
\file{app/config.py} reads the key's \emph{path} from the
\code{GOOGLE\_APPLICATION\_CREDENTIALS} environment variable and nothing else.
The key file lives outside the repo. \file{.env} is in \file{.gitignore}; only
\file{.env.example} is committed, containing placeholders such as
\code{GOOGLE\_APPLICATION\_CREDENTIALS=/absolute/path/to/your/service-account-key.json}.
I scanned the tracked files for private keys, API keys and tokens and found
none. This is the correct posture and it is worth stating in an interview.
\end{keyidea}

% =====================================================================
\section{The planner: the agentic step}

\file{app/planner.py} is 58 lines, and it is the reason the word ``agentic'' is
justified. The system prompt is blunt about the boundary:

\begin{lstlisting}[language=Python]
PLANNER_SYSTEM = (
    "You are a query planner for a trade-intelligence copilot. You decide which "
    "tools are needed to answer a question. You do not answer the question yourself."
)
\end{lstlisting}

The user prompt describes each of the three tools in one sentence, injects
\code{DATASET\_SCOPE}, and demands JSON back:

\begin{lstlisting}[language=Python,caption={The plan the model must return.}]
{
  "sql_needed": true/false,
  "sql_question": "<question restated for SQL>" or null,
  "retrieval_needed": true/false,
  "retrieval_query": "<search query>" or null,
  "hs_lookup": "<code or product name>" or null,
  "in_scope": true/false,
  "rationale": "<one sentence>"
}
\end{lstlisting}

Two design points deserve credit. First, the planner may set several fields at
once, which is what makes the multi-step questions work: compute a number with
SQL \emph{and} explain the rule with retrieval. Second, it \emph{rewrites} the
question per tool. \code{sql\_question} is the question restated for a
SQL-writing model and \code{retrieval\_query} is a search query, which is a
better input to an embedding model than a conversational sentence.

Finally, \code{plan()} normalizes the result, coercing every field with
\code{bool()} and \code{.get()} defaults, so a model that omits a key produces a
usable plan rather than a \code{KeyError}. That defensiveness is appropriate: an
LLM's JSON is untrusted input.

\begin{honest}{\code{in\_scope} is computed on every single query and never read}
The planner prompts for \code{in\_scope}, the model reasons about it, and
\code{plan()} carefully normalizes it with \code{result.get("in\_scope", True)}.
I grepped the entire codebase: \textbf{no other line ever reads it}. Refusal is
driven purely by the evidence check in \file{app/answer.py}. So this field costs
prompt and output tokens on every request and influences nothing.

It is not harmless. It reads like a safety control, and someone reviewing the
prompt would reasonably assume out-of-scope questions get short-circuited here.
They do not. Two honest defences: it is visible in the trace and the UI, where
it is genuinely useful for debugging; and it plausibly improves the plan by
making the model reason about scope before choosing tools, an effect similar to
asking for a rationale. But that is a hypothesis about prompting, not a
mechanism, and it is not what the field's name implies. Either wire it in (an
in-scope=false question could skip straight to refusal and save two model calls)
or delete it and keep the rationale.
\end{honest}

% =====================================================================
\section{The three tools}

\subsection{Tool 1: text-to-SQL}

\code{tools.text\_to\_sql()} assembles the schema DDL, the scope facts and the
SQL rules into one prompt, asks for \code{\{"sql": ..., "reason": ...\}}, and
then executes it through the guard. Its error handling is a good example of
distinguishing failure modes:

\begin{lstlisting}[language=Python]
try:
    result["rows"] = db.run_select(sql)
except db.SqlGuardError as e:
    result["error"] = f"guard_rejected: {e}"
except Exception as e:  # execution error (bad column, syntax)
    result["error"] = f"execution_error: {e}"
\end{lstlisting}

``The model tried something forbidden'' and ``the model wrote a broken query''
are different events and are recorded differently. Note also that the model is
permitted to answer \code{sql: null} when it judges the question unanswerable
from the tables, which becomes \code{no\_sql\_generated} and, downstream, no
evidence and a refusal.

\subsection{Tool 2 and 3, briefly}

\code{relevant\_chunks()} is a thin pair around the store: retrieve the top
\code{k}, then split into all hits and those clearing the threshold. Both halves
are returned, which is what lets the trace show a reader ``here is what came
back and here is what I threw away''.

\code{hs\_lookup()} is the only tool with no model in it at all: a
parameterised SQL join across the three hierarchy tables, matching a code prefix
(\code{term\%}) or a name substring (\code{\%term\%}), capped at 25 rows. Note
it uses \code{?} placeholders, so it is safe from SQL injection even though the
search term originates from an LLM.

\subsection{The SQL guard: the security boundary}
\label{sec:guard}

This is the most safety-critical code in the project, because an LLM is writing
executable statements against a database.

\begin{figure}[H]
\centering
\begin{tikzpicture}[node distance=7mm]
  \node[box] (in) {SQL text from the LLM};
  \node[dec, below=of in] (d1) {empty?};
  \node[dec, below=of d1] (d2) {contains \code{--}\\or \code{/*}?};
  \node[dec, below=of d2] (d3) {a \code{;} inside\\the body?};
  \node[dec, below=of d3] (d4) {starts with SELECT\\or WITH?};
  \node[dec, below=of d4] (d5) {a forbidden\\keyword?};
  \node[box, below=of d5, fill=warnred!12] (conn) {\textbf{Layer 2: the connection}\\\scriptsize mode=ro\&immutable=1\\\scriptsize PRAGMA query\_only = ON};
  \node[box, below=of conn, fill=warnred!12] (exec) {\textbf{Layer 3: the row cap}\\\scriptsize fetchmany(200)};
  \node[box, right=32mm of d3, fill=warnred!12] (err) {\textbf{SqlGuardError}\\\scriptsize nothing ever executes};

  \draw[flow] (in) -- (d1);
  \draw[flow] (d1) -- node[lbl,right]{no} (d2);
  \draw[flow] (d2) -- node[lbl,right]{no} (d3);
  \draw[flow] (d3) -- node[lbl,right]{no} (d4);
  \draw[flow] (d4) -- node[lbl,right]{yes} (d5);
  \draw[flow] (d5) -- node[lbl,right]{no} (conn);
  \draw[flow] (conn) -- node[lbl,right]{execute} (exec);
  \draw[flow] (d1) -| node[lbl,pos=0.22,above]{yes} (err);
  \draw[flow] (d2) -| node[lbl,pos=0.22,above]{yes} (err);
  \draw[flow] (d3) -- node[lbl,above]{yes} (err);
  \draw[flow] (d4) -| node[lbl,pos=0.22,above]{no} (err);
  \draw[flow] (d5) -| node[lbl,pos=0.22,above]{yes} (err);
\end{tikzpicture}
\caption{\code{validate\_select()} is layer 1: five gates, any of which aborts
before anything executes. Layers 2 and 3 assume layer 1 was bypassed anyway.}
\label{fig:guard}
\end{figure}

The layering is the actual insight, and \file{app/db.py} says so in a comment
(``Belt and braces''):

\begin{enumerate}
\item \textbf{The keyword and shape guard.} Reject comments (which can hide a
  payload), reject a semicolon anywhere but a single optional trailing one
  (which blocks \code{SELECT 1; DROP TABLE companies}), require the statement to
  begin \code{select} or \code{with}, and reject a regex of write and DDL
  keywords including \code{pragma}, \code{attach} and \code{vacuum}.
\item \textbf{The connection.} \code{sqlite3.connect("file:...?mode=ro\&immutable=1",
  uri=True)} opens the file read-only, and \code{PRAGMA query\_only = ON} tells
  SQLite itself to reject writes.
\item \textbf{The row cap.} \code{fetchmany(200)} means a \code{SELECT *} with
  no \code{LIMIT} cannot pull the whole table into memory or into a prompt.
\end{enumerate}

\begin{keyidea}{Why any one layer alone would be inadequate}
A keyword blocklist is a blocklist: someone will always think of a word you
forgot. A read-only connection would still allow a catastrophic
\code{SELECT} that exhausts memory. The row cap does nothing about writes. The
argument is not that each layer is strong; it is that a bypass must defeat all
three, and \file{tests/test\_sql\_guard.py} asserts each independently,
including that a raw \code{CREATE TABLE} on the connection raises
\code{sqlite3.OperationalError} even with the guard skipped entirely.
\end{keyidea}

\begin{honest}{The keyword regex reads string literals, and rejects valid queries}
\code{\_FORBIDDEN} is applied to the whole query text, with no awareness of SQL
string literals. Any forbidden word \emph{inside quotes} triggers it. I ran this
against the real guard and confirmed it:

\begin{lstlisting}[language=Python,caption={Verified against app/db.py. The first three are legitimate read-only queries.}]
"SELECT * FROM companies WHERE name LIKE '%Create%' LIMIT 5"       -> REJECTED
"SELECT * FROM companies WHERE name LIKE '%Update Traders%'"       -> REJECTED
"SELECT SUM(usd) FROM trade_transactions WHERE seller LIKE '%Delete Foods%'" -> REJECTED
"SELECT COUNT(*) FROM trade_transactions WHERE buyer LIKE '%Orvella%'"      -> accepted
\end{lstlisting}

A real Pakistani importer named ``Create Foods'' or ``Update Traders'' would be
permanently unqueryable, and the failure surfaces as
\code{guard\_rejected}, which is opaque to the user.

Now the fair defence, and it is a strong one: this fails \emph{safe}. A
false positive costs an answer; a false negative costs the database. Given the
choice, this is the right direction to err, and no company in the actual dataset
trips it. But it is a real limitation with a known fix (parse the SQL and inspect
only non-literal tokens, or use \code{sqlite3}'s \code{authorizer} callback,
which enforces at the engine level and cannot be fooled by quoting at all). Say
this in an interview before the interviewer says it to you.
\end{honest}

% =====================================================================
\section{Refusal, synthesis and citation}

\subsection{The deterministic refusal}

\file{app/answer.py} contains a nine-line function that carries a great deal of
weight:

\begin{lstlisting}[language=Python]
def deterministic_refusal(sql_result, chunks, hs_result) -> bool:
    """True when there is simply no evidence to answer from."""
    has_rows = bool(sql_result and sql_result.get("rows"))
    has_chunks = bool(chunks)
    has_hs = bool(hs_result and hs_result.get("matches"))
    return not (has_rows or has_chunks or has_hs)
\end{lstlisting}

If every tool came back empty, refuse now, before the synthesis call.

\begin{keyidea}{A cheap check beats an expensive one}
The alternative is to send ``here is no evidence, please refuse'' to a language
model and hope. This version cannot be argued out of it, cannot be
prompt-injected, costs zero tokens and zero seconds, and returns a fixed string.
It is faster, cheaper and more reliable, all three at once. That is a rare
trade in engineering and it is worth naming.
\end{keyidea}

\subsection{How refusal actually happened in the real run}

I examined all six unanswerable items in \file{eval/transcripts.json} to see
which mechanism fired. All six refused, giving the published 1.00 precision and
1.00 recall. Every one took the deterministic path, and every trace was exactly
two steps: \code{['plan', 'refusal']}.

\begin{honest}{``Deterministic'' refusal is downstream of an LLM's judgment}
Look closely at that trace. \code{['plan', 'refusal']} means \textbf{no tool ran
at all}. The deterministic guard did not sift through evidence and find it
wanting; it fired because the planner returned \code{sql\_needed: false},
\code{retrieval\_needed: false}, \code{hs\_lookup: null} and gathered nothing.
(The planner also set \code{in\_scope: false} on all six, correctly, and as
established, nothing read it.)

So the refusal is deterministic in its \emph{final step} and LLM-driven in its
\emph{decisive} one. The guard is a reliable backstop for ``the tools found
nothing''; it is not what recognised the question as out of scope. The planner
did that, and the planner is a language model.

This matters concretely. Consider \code{refuse-06}, ``What was the total value
of molasses exports?''. Molasses (HS 1703) \emph{is} in the corpus and in the
HS hierarchy. Had the planner set \code{hs\_lookup: "molasses"}, that tool would
have returned real matching rows, \code{has\_hs} would be true, the
deterministic guard would \emph{not} have fired, and the refusal would have
rested entirely on the synthesizer noticing that rows about molasses imports
cannot answer a question about molasses exports. It did not come to that, and
6/6 is a real result. But the mechanism has an LLM on its critical path, and
``enforced deterministically'' oversells it slightly. The correct claim is:
``the refusal is emitted by a deterministic check, which fires when the planner
selects no tools; on the eval set the planner selected no tools every time.''
\end{honest}

\subsection{Synthesis}

\code{synthesize()} builds one evidence block containing the SQL query, any SQL
error, the rows as JSON (capped at 50), the HS-lookup matches (capped at 25) and
the regulation chunks with their ids. It then instructs:
\emph{Answer ONLY from the evidence above. Do not use outside knowledge for
facts, numbers, or specific rules.} And it demands JSON with \code{supported},
\code{answer}, \code{cited\_chunks} and \code{used\_sql}.

Note the belt-and-braces on the way out: if the model says
\code{supported: false} but writes no text, the code substitutes
\code{REFUSAL\_TEXT}. The system never returns an empty answer.

\subsection{Citations, and one latent problem}

\file{app/agent.py} assembles the citation block: the SQL, up to 20 rows, up to
25 HS matches, and the full text of every chunk the model actually cited
(filtered by an intersection with \code{cited\_chunks}). This block is what the
UI renders and, crucially, what the eval's judge is later shown.

\begin{honest}{The judge's view of the evidence depends on the model's self-report}
One line in \file{app/agent.py} deserves scrutiny:

\begin{lstlisting}[language=Python]
"sql_rows": (sql_result.get("rows") or [])[:20] if (sql_result and result["used_sql"]) else [],
\end{lstlisting}

\code{used\_sql} is not observed by the code. It is a field the \emph{synthesizer
model} returned about itself. If the model uses a number from the SQL rows but
reports \code{used\_sql: false}, the rows are dropped from the citations, and
the judge, which reconstructs evidence from the citations, then sees an answer
containing a number with no source for it, and correctly calls it unfaithful.
The answerer would be right and the harness would score it wrong.

This is exactly the class of bug the project already found and fixed once: the
README records that HS-lookup rows were originally missing from the citations,
so the judge flagged good answers as unfaithful until they were added. The same
structural weakness remains for \code{sql\_rows}, gated on a self-report rather
than on whether rows existed.

Credit where due: I checked all 28 transcripts for a case where SQL returned
rows but \code{sql\_rows} was empty in the citations, and found \textbf{zero}.
The bug did not bite in the published run (15 SQL steps returned rows, and all
15 items cite them), so the 95.5\% figure is not contaminated. It is a latent fragility, not an active error, and the fix is one
line: gate on \code{sql\_result.get("rows")} rather than on
\code{result["used\_sql"]}.
\end{honest}

% =====================================================================
\section{Cost accounting, the server and the UI}

\file{app/llm.py} wraps generation and meters it. A \code{Usage} record holds a
model name and input and output token counts, and prices itself against the
table in \file{app/config.py}. A \code{CostMeter} accumulates \code{Usage}
records across the three or four calls a single question makes, and one meter is
threaded through the whole request, which is why the API can report a genuine
per-query cost rather than an estimate.

\begin{honest}{The prices are hardcoded list rates, not billed amounts}
\code{PRICING} in \file{app/config.py} holds figures such as \$0.30 per million
input tokens for flash. The token \emph{counts} are real, taken from the API's
own \code{usage\_metadata}. The dollar figures are published list prices copied
into a dict, overridable by environment variable, and the comment says
``Approximate published rates''. So \$0.00093 per query is a real token count
multiplied by an assumed price. It is not read from a billing API and will
silently drift if Google changes pricing. This is a reasonable engineering
choice, honestly commented, and worth stating precisely rather than saying
``the system tracks cost''.
\end{honest}

\code{\_parse\_json} is small and pragmatic: it strips Markdown code fences and,
if the model wrapped its JSON in chatter, trims to the outermost braces. Ugly,
necessary, and exactly what everyone who has shipped an LLM feature writes.

\file{app/server.py} is 40 lines: \code{GET /health} reports whether the
database and vector files are present and whether credentials are configured
(without revealing them), \code{POST /ask} calls the agent, and \code{GET /}
returns the UI.

\begin{jargon}{FastAPI}
A Python web framework. A \textbf{web framework} handles the plumbing of
receiving HTTP requests and returning responses, so the author writes only the
logic. The \code{AskRequest} class declares that \code{/ask} expects
\code{\{"question": "..."\}}, and FastAPI rejects malformed requests before any
code runs.
\end{jargon}

The UI is a single self-contained HTML file with no build step and no
dependencies. It renders the answer with a grounded or refused badge, the SQL
executed, the rows, the HS matches, the cited chunk text, the full tool-call
trace, and the token count, cost and latency.

\begin{keyidea}{The UI shows the trace, and that is a statement of intent}
Most demos show the answer. This one shows the plan, every tool call, what
retrieval returned \emph{and what it discarded}, the exact SQL, and the cost. A
user can audit any answer without reading the code. Making the machinery visible
is the same instinct as publishing the eval.
\end{keyidea}

\begin{honest}{\code{jinja2} is a dependency that is never used}
\file{requirements.txt} pins \code{jinja2==3.1.6} under the comment ``Config +
templating''. \file{app/server.py} reads the template with
\code{Path(...).read\_text()} and does no templating at all. I grepped every
Python file: jinja2 is never imported. It is a leftover from an approach that
was abandoned, and the file even lives in a folder called \file{templates/}
which is never templated. Harmless, but it is dead weight in a dependency list
and the sort of thing worth removing before someone asks about it.
\end{honest}

% =====================================================================
\section{The evaluation layer}

This is the part of the project worth the most in an interview.

\subsection{The test set}

\file{eval/testset.jsonl} holds 28 items, one JSON object per line.

\begin{jargon}{JSONL}
JSON Lines: one complete JSON object per line of the file. Handy for test sets
because you can append an item, or read them one at a time, without parsing the
whole file.
\end{jargon}

\begin{table}[H]
\centering
\small
\begin{tabular}{@{}llll@{}}
\toprule
\textbf{Group} & \textbf{n} & \textbf{Gold answer is} & \textbf{Scored on} \\
\midrule
\code{sql} & 12 & a number (with tolerance) or a substring & SQL correctness \\
\code{retrieval} & 7 & the document(s) that should be found & retrieval hit rate \\
\code{multi} & 3 & both a number and a document & both \\
\code{refusal} & 6 & nothing; it must refuse & refusal precision/recall \\
\bottomrule
\end{tabular}
\caption{The 28 items. Every item is also judged for faithfulness unless it was
supposed to be refused. SQL correctness is scored over \code{sql} plus
\code{multi}, which is $12 + 3 = 15$: that is where ``15/15'' comes from.}
\end{table}

The tolerance design is careful: \code{rel\_tol: 0.0} for integer counts
(a count is exactly right or it is wrong) and \code{rel\_tol: 0.01} for money,
absorbing rounding in how the model formats a dollar figure.

\subsection{The harness}

\begin{figure}[H]
\centering
\resizebox{\textwidth}{!}{
\begin{tikzpicture}[node distance=7mm]
  \node[box] (ts) {\file{testset.jsonl}\\\scriptsize 28 items + gold};
  \node[box, right=12mm of ts, fill=accent!12] (ask) {\textbf{agent.ask()}\\\scriptsize the real copilot,\\\scriptsize live API};
  \node[box, right=12mm of ask, fill=accent!12] (judge) {\textbf{judge}\\\scriptsize gemini-2.5-pro\\\scriptsize faithful?};
  \node[store, right=12mm of judge] (tr) {\file{transcripts.json}\\\scriptsize answer + trace\\\scriptsize + cost + verdict};

  \node[box, below=20mm of ts, fill=warnred!12] (metrics) {\textbf{compute\_metrics()}\\\scriptsize OFFLINE, no API};
  \node[store, left=10mm of metrics] (hl) {\file{human\_labels.json}\\\scriptsize 12 independent labels};

  \node[box, right=14mm of metrics] (m1) {SQL correctness\\\scriptsize number/substring vs gold};
  \node[box, right=12mm of m1] (m2) {retrieval hit rate\\\scriptsize gold doc retrieved?};
  \node[box, below=6mm of m1] (m3) {refusal matrix\\\scriptsize tp/fp/tn/fn};
  \node[box, below=6mm of m2] (m4) {faithfulness\\\scriptsize + judge vs labels};
  \node[store, right=12mm of m2, yshift=-9mm] (res) {\file{results.json}\\\scriptsize the published numbers};

  \draw[flow] (ts) -- (ask); \draw[flow] (ask) -- (judge); \draw[flow] (judge) -- (tr);
  \draw[dflow] (tr.south) -- ++(0,-7mm) -| node[lbl,pos=0.25,above]{saved once, scored forever} (metrics.north);
  \draw[flow] (hl) -- (metrics);
  \draw[flow] (metrics) -- (m1); \draw[flow] (metrics) -- (m3);
  \draw[flow] (m1) -- (m2); \draw[flow] (m3) -- (m4);
  \draw[flow] (m2) -- (res); \draw[flow] (m4) -- (res);
\end{tikzpicture}}
\caption{The harness, split in two on purpose. \code{--run} (top) is the
expensive half and happens once. \code{--metrics} (bottom) recomputes every
published number from the saved transcripts and costs nothing.}
\label{fig:eval}
\end{figure}

\begin{keyidea}{The two-phase split is the best structural decision in the project}
Because the transcripts are saved, the scoring can be re-run, corrected,
extended and \emph{independently checked} without spending another cent on the
API. It also means the published results are auditable: anyone can recompute
them from the committed transcripts. I did exactly that
(Section~\ref{sec:verify}), and it is only possible because of this split.
\end{keyidea}

\subsection{How each metric is computed}

\begin{itemize}
\item \textbf{SQL correctness.} \code{extract\_numbers()} regex-scrapes every
  number out of the answer \emph{text}, strips commas, and compares each against
  gold within tolerance; or, for a \code{contains} gold, does a lowercase
  substring test. Note this scores the \textbf{whole pipeline}: planner, SQL
  generation, guard, execution and synthesis all had to work for the right
  number to appear in the prose.
\item \textbf{Retrieval hit.} Find the \code{retrieval} step in the trace, take
  the chunk ids that came back, strip each at the \code{\#} to get a document
  name, and ask whether any gold document is among them. Note it uses
  \emph{all} hits, not just those above the threshold.
\item \textbf{Refusal.} A confusion matrix over all 28 items, comparing
  ``refused'' (\code{supported == false}) against \code{expect\_refusal}.
\item \textbf{Faithfulness.} The judge's verdict, counted only over items that
  were answered and were not supposed to be refused (22 of the 28).
\item \textbf{Judge vs labels.} For each of the 12 labelled ids, does the
  author's label match the judge's verdict?
\end{itemize}

\begin{jargon}{Precision and recall}
Two ways of being wrong. Here, \textbf{recall} = of the questions that should
have been refused, how many were? (Did it catch them all?) \textbf{Precision} =
of the times it refused, how often was refusing correct? (Is it refusing things
it could have answered?) A system that refuses everything has perfect recall and
terrible precision. You need both, which is why both are reported.
\end{jargon}

\subsection{The published results}

\begin{table}[H]
\centering
\small
\begin{tabular}{@{}lll@{}}
\toprule
\textbf{Metric} & \textbf{Result} & \textbf{My verification status} \\
\midrule
SQL answer correctness & 15 / 15 = 100\% & recomputed offline; gold re-derived from DB \\
Retrieval hit rate & 8 / 10 = 80\% & recomputed offline; cause confirmed \\
Refusal precision / recall & 1.00 / 1.00 & recomputed offline; mechanism inspected \\
Judge faithfulness & 21 / 22 = 95.5\% & recomputed offline; the miss read \\
Judge vs independent labels & 11 / 12 = 91.7\% & recomputed offline; each label read \\
Mean cost per query & \$0.00093 & recomputed; rests on list-price table \\
Total cost of the run & \$0.0260 & recomputed \\
Latency p50 / p95 & 9.0 s / 27.0 s & recomputed; 3 items include 429 waits \\
\bottomrule
\end{tabular}
\caption{The README's headline table, with what I could and could not confirm.
Section~\ref{sec:verify} is precise about the difference.}
\end{table}

% =====================================================================
\section{What I verified by actually running it}
\label{sec:verify}

Claims in a README are assertions. Here is what I executed, and what it proved.

\begin{keyidea}{The live run was re-done when the data was replaced}
On 2026-10-14 the synthetic database replaced the original data, and
\code{python -m eval.run\_eval --run} was run once against the live API (28
items, about \$0.026), along with the four live tests. Everything else below
runs offline.
\end{keyidea}

\subsection{The gold answers are genuinely ground truth}

The strongest claim in the project is that the gold numbers come from the
database rather than from anybody's judgment. They are written by
\file{eval/build\_testset.py}, which runs a query for each one, and I re-derived
them with independent queries against the shipped \file{data/trade.sqlite}.
\textbf{They reproduce exactly.} For example: total USD 274,995,234.91; row count
2,000; China 208 shipments; the top country by value is Brazil at
112,021,544.46; Orvella Foods 191 shipments; HS 1702.1100 227 shipments; HS 1704
total 11,665,069.33.

I also confirmed every scope fact asserted in \code{DATASET\_SCOPE}:
\code{trade\_type} has exactly one distinct value, \code{'import'};
\code{reporting\_date} runs 2025-01-01 to 2025-12-31; there are 28 distinct
countries and 28 companies. Two builds of the database are byte-identical.

\subsection{The metrics reproduce exactly}

I created a clean virtual environment, installed the dependencies, and ran
\code{python -m eval.run\_eval --metrics} with credentials explicitly unset, so
no API call was possible. It recomputed every metric from the committed
transcripts and produced output \textbf{byte-identical} to the committed
\file{eval/results.json}: 15/15 SQL, 8/10 retrieval, refusal 6 tp / 0 fp / 22 tn
/ 0 fn, faithfulness 21/22, agreement 11/12, cost \$0.025975, p50
8,991 ms.

\subsection{The test suite}

\code{pytest} with no credentials: \textbf{29 passed, 4 skipped}. That matches
the README's claim exactly (``33 tests total: 29 run offline ... 4 are marked
\code{live}''). The live tests skipped themselves cleanly via the
\code{conftest.py} hook rather than failing, and with a Vertex key set all four
pass against the API.

\subsection{So is ``15/15'' verified?}

\begin{honest}{Precisely what the 15/15 rests on}
All three links in the chain are verified for the current numbers:

\begin{enumerate}
\item \textbf{The gold answers are correct.} Verified independently against the
  database. All 13 numbers exact.
\item \textbf{The scoring of the transcripts yields 15/15.} Verified by
  re-running the scorer offline; byte-identical output.
\item \textbf{The transcripts are the output of a real live run against Gemini.}
  Verified: they were produced by the 2026-10-14 run described above.
\end{enumerate}

That run hit Vertex rate limits. Three items (\code{sql-10}, \code{sql-11},
\code{ret-07}) waited out \code{429 RESOURCE\_EXHAUSTED} responses, so their
recorded latencies (27, 46 and 30 seconds) include waiting. That is why the p95
is reported twice: 27.0 s as measured, 12.5 s without those three. Nothing was
re-run to make a number look better; the two retrieval misses and the one judged
unfaithful answer are all still in the transcripts.
\end{honest}

\subsection{What the two retrieval misses actually were}

The README says the 2 misses were questions answered from the HS-lookup tool
instead of retrieval. I checked the traces of all 10 retrieval-scored items.
\code{ret-07} and \code{multi-01} have \textbf{no retrieval step in the trace at
all}: the planner never invoked retrieval, so no chunk could be retrieved, so
the gold document could not be hit. The other 8 all retrieved their gold
document. The README's account is accurate.

\begin{honest}{The retrieval metric conflates two different failures}
\code{retrieval\_hit()} returns \code{False} both when retrieval ran and fetched
the wrong document, and when retrieval never ran. Those are different bugs with
different fixes: a bad embedding versus a bad routing decision. Both of this
project's misses are the second kind, which means \textbf{80\% is a planner
score, not a retrieval score}. The retrieval component did not miss a single
document it was actually asked to find; it went 8 for 8. The metric's name
points at the wrong component, and the fix the README proposes (make the planner
prefer retrieval for definitional questions) is aimed correctly even though the
metric is not.
\end{honest}

\subsection{The one unfaithful verdict, read in full}

The single unfaithful verdict is \code{sql-12}, ``How many shipments were
recorded under HS code 1702.1100 (lactose)?''. The answer gave the right count,
227, and repeated the question's own label. The judge ruled unfaithful, with the
recorded reason: \emph{``The evidence supports the number of shipments (227) for
the given HS code, but it does not provide the description `(lactose)' for that
code. This description is an invented detail not found in the evidence.''}

\begin{keyidea}{Both readings are defensible, and that is the point}
The judge is literally correct: the evidence did not establish that mapping.
The independent labels call it faithful: the load-bearing number was right and the
label was user-supplied. It is the single disagreement in the 11/12 judge check, as
the same kind of answer was on the first version. Publishing the disagreement,
rather than relabelling until agreement hits 12/12, is what makes 91.7\% mean
something.
\end{keyidea}

\begin{keyidea}{Enumeration is where grounded models leak}
In the first version, the other two unfaithful verdicts were the same failure:
asked to enumerate what an HS heading covers, the model padded its list with
plausible product names that were not in the evidence rows. Point-fact answers
stayed grounded; ``list everything under this heading'' is where the model
reached past its evidence, because a short list \emph{feels} incomplete. In the
synthetic re-run both of those questions (\code{ret-07}, \code{multi-01}) were
judged faithful, most likely because the synthetic HS hierarchy is small enough
for the lookup to return every item. That is luck, not a fix. The fix is
structural, not a stronger instruction: post-filter enumerated items against the
evidence rows.
\end{keyidea}

% =====================================================================
\section{Consolidated honest findings}
\label{sec:findings}

Everything below is either verified by me or read directly out of the source.
The first three are the project's own, already documented in its README, and
they are here because they are good.

\subsection{Findings the project already owns, and gets right}

\begin{enumerate}
\item \textbf{The wildcard bug.} In the first version, text-to-SQL matched one
  importer's full company name with no wildcards, matched nothing, and reported
  0 shipments for an importer that had dozens. The answer was fluent and completely wrong. The gold number
  caught it instantly; the fix was a sharper instruction in the schema prompt.
  This is the single best story in the project, because it is a concrete
  demonstration that the eval earns its keep.
\item \textbf{Faithful is not correct.} That zero-row answer was judged
  \emph{faithful}: it accurately reported what its query returned. This is why
  grounding and ground-truth correctness are measured separately. Anyone
  reporting a single ``accuracy'' number for a RAG system has not thought this
  through.
\item \textbf{The judge was blind to evidence the answerer used.} The first run
  flagged HS-lookup answers as unfaithful because the judge could not see the
  HS rows. Fixing the plumbing turned the remaining flags into genuine findings.
  A harness bug that manufactures failures is as dangerous as one that hides
  them.
\end{enumerate}

\subsection{Findings from this review}

\begin{enumerate}
\item \textbf{\code{in\_scope} is dead.} Computed on every query, read by
  nothing. Reads like a safety control, is not one.
\item \textbf{The SQL guard rejects valid queries whose string literals contain
  a keyword.} Verified by execution. Fails safe, but a company named ``Create
  Foods'' is unqueryable.
\item \textbf{\code{sql\_rows} citation is gated on the model's self-reported
  \code{used\_sql}.} A latent repeat of a bug already fixed once elsewhere.
  Verified not to have affected the published run (zero instances).
\item \textbf{``Deterministic refusal'' is deterministic in its last step
  only.} The decisive choice is the planner's, and the planner is an LLM.
\item \textbf{The retrieval hit rate is really a planner score.} Both misses were
  routing failures; retrieval went 8 for 8 on what it was asked to find.
\item \textbf{\code{jinja2} is an unused dependency.}
\item \textbf{Costs are list prices in a dict, not billed amounts.} Token counts
  are real; the dollar conversion is an assumption that will drift.
\end{enumerate}

\subsection{Limitations the project states, and should keep stating}

\begin{itemize}
\item \textbf{One run, 28 items, temperature 0.} Every number here is a single
  sample. Temperature 0 makes the models near-deterministic but not guaranteed
  identical, and 15/15 on 15 questions is a good result with wide error bars.
  The honest reading is ``the paths work'', not ``it is 100\% accurate''.
\item \textbf{SQL correctness is scraped from prose.} Any number within
  tolerance anywhere in the answer scores a hit, so a wrong answer containing
  the right number would pass. \file{docs/EVALUATION.md} states this itself.
\item \textbf{The judge and the answerer are the same model family.} Shared
  blind spots are possible; the independent labels only partly cover this. A judge
  from a different family would be a stronger cross-check.
\item \textbf{The corpus is authored summaries, not primary law.}
\item \textbf{Latency is high (p50 about 9 s)} because three or four model calls
  run in sequence. Merging planning with text-to-SQL, and running independent
  tools concurrently, would cut it materially. The code runs the tools strictly
  one after another even when they are independent.
\end{itemize}

% =====================================================================
\section{Defending this project in an interview}

\begin{keyidea}{The five things to say}
\begin{enumerate}
\item \textbf{Lead with the eval, not the copilot.} Anyone can wire an LLM to a
  database. The differentiator is 28 items with gold answers computed from the
  database, published metrics including the failures, and a judge checked
  against independent labels.
\item \textbf{Tell the wildcard story.} A fluent answer said 0 shipments; the
  truth was dozens; the gold number caught it in one run. That is the whole argument
  for evaluation in thirty seconds.
\item \textbf{Separate faithful from correct}, and explain that the same bug
  proved why: the wrong answer was perfectly faithful to its own wrong query.
\item \textbf{Defend the hand-built vector store.} 48 chunks, exact brute-force
  search, one matrix-vector product, no server. pgvector was unavailable; for
  this corpus the simpler thing is also the better thing.
\item \textbf{Volunteer the limits.} 80\% retrieval is really a routing problem.
  Enumeration is where grounding leaks. \code{in\_scope} is dead code. Say these
  before you are asked.
\end{enumerate}
\end{keyidea}

\subsection{Questions you should expect}

\begin{itemize}
\item \emph{``What stops the model dropping your tables?''} Three layers:
  keyword and shape guard, a connection opened \code{mode=ro\&immutable=1} with
  \code{PRAGMA query\_only}, and a row cap. A test asserts a write fails even
  with the guard bypassed. Then volunteer the string-literal false positive.
\item \emph{``How do you know it refuses instead of guessing?''} Six unanswerable
  items, all refused, zero false refusals. Then be precise: the deterministic
  check emits the refusal, but it fires because the planner chose no tools.
\item \emph{``Isn't using an LLM to grade an LLM circular?''} Yes, which is why
  the judge is checked against 12 independent labels: 11/12 on both the first
  version and the synthetic re-run, labelled afresh each time, with the one
  disagreement published.
\item \emph{``Faithfulness is 95.5\%; why should I trust it?''} Because the judge
  agrees with independent labels 11 times in 12, and its one miss is a judgment
  call about a label the user supplied. On the first version it was 86\%, and two
  of the three misses were the model padding an enumeration beyond its evidence, a
  real failure mode with a structural fix.
\item \emph{``Would this scale?''} The vector store would need a real ANN index
  somewhere around six figures of chunks. The bigger constraint is latency:
  three or four sequential model calls, about 9 seconds. Merge the planner with
  text-to-SQL and parallelise independent tools.
\end{itemize}

% =====================================================================
\section{Glossary}

\begin{description}[leftmargin=!,labelwidth=3.4cm,style=nextline]
\item[Agentic retrieval] Retrieval where a model first decides which sources to
  consult rather than always running one fixed search.
\item[Bearer token] A short-lived string proving a request is authorised,
  minted here from a service-account key.
\item[Chunk] One retrievable piece of a document; here, one \code{\#\#} section.
\item[Cosine similarity] How closely two vectors point the same way; 1.0 is
  identical direction. With normalized vectors it is just the dot product.
\item[Embedding] A list of numbers (768 here) representing the meaning of text.
\item[Faithfulness] Whether an answer is supported by the evidence it was given.
  Not the same as correctness.
\item[Ground truth] The known-correct answer, established independently. Here,
  computed by querying the database directly.
\item[Hallucination] Fluent, plausible, false model output.
\item[HS code] International product classification. This dataset is chapter 17,
  sugar.
\item[JSONL] One JSON object per line.
\item[LLM-as-judge] Using a model to grade another model's output.
\item[Precision / recall] Of the refusals, how many were right / of the things
  needing refusal, how many were caught.
\item[RAG] Retrieval-augmented generation: fetch the facts, then let the model
  write using only those facts.
\item[SQL guard] The code that refuses to execute anything but a single, safe
  \code{SELECT}.
\item[SQLite] A complete relational database that is a single file, no server.
\item[Text-to-SQL] Having a model translate a question into a database query and
  running it for real.
\item[Vector store] The structure holding embeddings and searching them; here a
  48x768 numpy matrix.
\end{description}

\vfill
\begin{center}
\small\itshape
Built from the source at \file{trade-intelligence-copilot}. Every metric quoted
was recomputed offline from the committed transcripts, and every gold answer
re-derived from the shipped database, before it was written down.
\end{center}

\end{document}
